%  LaTeX support: latex@mdpi.com 
%  For support, please attach all files needed for compiling as well as the log file, and specify your operating system, LaTeX version, and LaTeX editor.

%=================================================================
\documentclass[jmse,article,submit,pdftex,moreauthors]{Definitions/mdpi} 

%=================================================================

% MDPI internal commands
\firstpage{1} 
\makeatletter 
\setcounter{page}{\@firstpage} 
\makeatother
\pubvolume{1}
\issuenum{1}
\articlenumber{0}
\pubyear{2024}
\copyrightyear{2024}
%\externaleditor{Academic Editor: Firstname Lastname}
\datereceived{27 June 2024} 
\daterevised{18 July 2024} 
\dateaccepted{} 
\datepublished{} 
\hreflink{https://doi.org/} % If needed use \linebreak
%\doinum{}

%=================================================================
% Add packages and commands here. The following packages are loaded in our class file: fontenc, inputenc, calc, indentfirst, fancyhdr, graphicx, epstopdf, lastpage, ifthen, lineno, float, amsmath, setspace, enumitem, mathpazo, booktabs, titlesec, etoolbox, tabto, xcolor, soul, multirow, microtype, tikz, totcount, changepage, attrib, upgreek, cleveref, amsthm, hyphenat, natbib, hyperref, footmisc, url, geometry, newfloat, caption

\graphicspath{{figures/}} % path to folder with figures
\usepackage{subcaption}
\usepackage{listings}
\usepackage{xcolor}
\usepackage{gensymb} % for \textdegree
\usepackage{tabularx}
\newcolumntype{Y}{>{\centering\arraybackslash}X}
\setlength{\extrarowheight}{-20pt}
\renewcommand{\arraystretch}{0.7}
\usepackage{amsmath,mathtools,amsthm}
	\setcounter{MaxMatrixCols}{15}
\usepackage{array}
\usepackage{amssymb}
\usepackage{amsfonts}
\usepackage{float}
\usepackage{chngcntr}
\counterwithin*{equation}{section}
\usepackage[super]{nth}
\usepackage{longtable}
\usepackage{tabularx}
\usepackage{adjustbox}
\usepackage{graphicx}%
\usepackage{multirow}%
\usepackage{amsmath,amssymb,amsfonts}%
\usepackage{amsthm}%
\usepackage{mathrsfs}%
\usepackage[title]{appendix}%
\usepackage{textcomp}%
\usepackage{manyfoot}%
\usepackage{booktabs}%
\usepackage{algorithm}%
\usepackage{algorithmicx}%
\usepackage{algpseudocode}%
\usepackage{listings}%
\usepackage{rotating}
\usepackage{makecell}


\definecolor{deepblue}{rgb}{0,0,0.5}
\definecolor{deepred}{rgb}{0.6,0,0}
\definecolor{deepmagenta}{RGB}{195,12,214}
\definecolor{deepgreen}{RGB}{29,122,30}
\definecolor{mintgreen}{RGB}{192,249,192}
\definecolor{codepurple}{RGB}{204, 0, 153}
\definecolor{LightCyan}{rgb}{0.88,1,1}
\definecolor{GainsboroPurple}{RGB}{247,176,247}
\definecolor{LightSlateGrey}{RGB}{124,118,118}
%    
\lstdefinestyle{mystyle}{
    commentstyle=\color{LightSlateGrey},
    keywordstyle=\color{deepgreen}\bfseries,
    emphstyle=\color{deepred},
    numberstyle=\tiny\color{deepblue},
    stringstyle=\color{codepurple},
    basicstyle=\ttfamily\scriptsize,
    breakatwhitespace=false,         
    breaklines=true,                 
    captionpos=b,                    
    keepspaces=true,                 
    numbers=left,                    
    numbersep=5pt,                  
    showspaces=false,                
    showstringspaces=false,
    showtabs=false,                  
    tabsize=2
}
\lstset{style=mystyle}

%=================================================================
% Full title of the paper (Capitalized)
\Title{Artificial Intelligence for Computational Remote Sensing: Quantifying Patterns of Land Cover Types Around Cheetham Wetlands, Port Phillip Bay, Australia}

% MDPI internal command: Title for citation in the left column
\TitleCitation{Artificial Intelligence for Computational Remote Sensing: Quantifying Patterns of Land Cover Types Around Cheetham Wetlands, Port Phillip Bay, Australia}

% Author Orchid ID: enter ID or remove command
\newcommand{\orcidauthorA}{0000-0002-5759-1089} % Polina Add \orcidA{} behind the author's name

% Authors, for the paper (add full first names)
\Author{Polina Lemenkova $^{1}$,$^{2}$*\orcidA{}}

%\longauthorlist{yes}

% MDPI internal command: Authors, for metadata in PDF
\AuthorNames{Polina Lemenkova}

% MDPI internal command: Authors, for citation in the left column
\AuthorCitation{Lemenkova, P.}
% If this is a Chicago style journal: Lastname, Firstname, Firstname Lastname, and Firstname Lastname.

% Affiliations / Addresses (Add [1] after \address if there is only one affiliation.)
\address{%
$^{1}$ \quad Department of Geoinformatics, Faculty of Digital and Analytical Sciences, Paris Lodron Universität Salzburg, Schillerstraße 30, A-5020 Salzburg, Austria.

$^{2}$ \quad Department of Biological, Geological and Environmental Sciences, Alma Mater Studiorum -- Università di Bologna, Via Irnerio 42, Bologna, IT-40126, Emilia-Romagna, Italy.
}

% Contact information of the corresponding author
\corres{Correspondence: polina.lemenkova@plus.ac.at; polina.lemenkova2@unibo.it; Tel.: (+43)-677-6173-2772 or (+39)-344-69-28-732}

% Current address and/or shared authorship
%\firstnote{Current address: Affiliation 3.} 

% Abstract (Do not insert blank lines, i.e. \\) 
\abstract{This paper evaluates the potential of using Artificial Intelligence (AI) and Machine Learning (ML) \textcolor{blue}{approaches} for classification of \textcolor{blue}{the} Landsat satellite imagery \textcolor{blue}{for environmental coastal mapping. The aim is} to identify changes in patterns of land cover types in a coastal area around Cheetham Wetlands, Port Phillip Bay, Australia. The scripting approach of Geographic Resources Analysis Support System (GRASS) Geographic Information System (GIS) uses AI-based methods of image analysis to accurately discriminate land cover types. Four ML algorithms were applied\textcolor{blue}{, tested and compared} for supervised classification\textcolor{blue}{. Technical approaches are based on} using "r.learn.train" module that employs the Scikit-Learn library of Python. The \textcolor{blue}{methodology includes following algorithms}: 1) Random Forest (RF), 2) Support Vector Machine (SVM), 3) ANN-based approach using Multi-layer Perceptron (MLP) Classifier; and 4) Decision Tree Classifier (DTC). The tested methods of AI demonstrated \textcolor{blue}{robust} results of image classification with the highest overall accuracy that exceeds 98\% reached by the SVM and RF scenarios. The presented scripting approach of GRASS GIS accurately detected changes in land cover types in southern Victoria over the period of 2013-2024. \textcolor{blue}{From our findings, the use AI and ML algorithms offers effective solutions for coastal monitoring by analysis of change detection using multi-temporal RS data.} The demonstrated \textcolor{blue}{methods have} potential applications in coastal and wetland monitoring, environmental analysis and urban planning based on Earth observation data.}

% Keywords
\keyword{image processing; Earth observation; coasts; GIS; cartography; geoinformatics; landscape analysis; environmental monitoring; sensors; machine learning; Melbourne; Victoria} 

% The fields PACS, MSC, and JEL may be left empty or commented out if not applicable
%\PACS{J0101}
%\MSC{}
%\JEL{}

\begin{document}

%----------- section ----------->
\section{Introduction}

Coastal regions are heterogeneous areas with\hl{ }complex landscape patterns that reflect the interplay between natural ecosystems and human activities. The environmental settings of such areas reflect \hl{tight} links between the terrestrial and marine natural ecosystems. On the other hand, coastal lands experience high anthropogenic pressure such as intensive urban construction, industrial and agriculture activities, marine transport systems and harbours. Monitoring the dynamics of the shorelines is essential for effective land management, natural resource planning, and urban development. Remote sensing (RS) data, such as satellite images, present a robust, reliable and cost-efficient information source to monitor coastal landscapes \cite{AGATE2024108639,Lemenkova202405,LIU2024120096}. 

Important advantage of the RS data consists in the long and ongoing tradition of \hl{their} using\hl{ }in environmental studies \hl{in general and} of coastal regions\hl{ in particular}. This became possible due to the rich collection of RS data available through different satellite \hl{missions}. As a result, a systematic and comprehensive observation of the Earth's surface made through the satellite images provided \hl{ }rich informational background for \hl{deep} insights into the environmental analysis of climate-related processes. For instance, time series of satellite images are useful for analysis and modelling of landscape dynamics. Comparison of such data enables to reveal \hl{important ongoing} trends in coastal processes \hl{of} the\hl{ }marine \hl{ecosystems, for instance}, to visualise flood extents \cite{10282975}, to highlight the intensity of hazards \cite{8517662} \hl{or} to evaluate the potential risks\hl{ }\cite{8900280}. 

RS data are widely used for monitoring coastal and marine environments. \hl{Examples of the use of the RS data in environmental studies} include\hl{, for instance, the analysis of }coastal erosion \cite{MCCARROLL2024107146}, monitoring vegetation and landscape dynamics \cite{Munroe2022,Morgan}, \hl{detection of} seagrass distribution \cite{Traganos}, assessment of water pollution and\hl{ }eutrophication \cite{ZHU2024262}, visualization of \hl{the }urban sprawl in coastal areas \cite{su14095700}, mapping land cover changes \cite{Lemenkova202313}, evaluating water quality \cite{ZHU2022116187,Lemenkova202322,Kim}, computing eutrophic sea level rise or shoreline extraction \cite{MCALLISTER2022104102}\hl{. }Among these \hl{examples, detecting}, land cover \hl{changes} plays a special role for coastal monitoring. \hl{In fact,} it enables to evaluate environmental trends detect dynamics and provide data for prognosis of possible scenario development. Besides forecasting, detecting changes is crucial for evaluating climate-environmental interactions, \hl{as} changed environment affects human well-being and destructs social systems through worsen living conditions. With this regards,\hl{ }coastal monitoring enables to better understand the complexity of such processes which is essential for environmental planning and policy making \cite{MURRAY2000357,COUTTS2016637}.

Many studies have been undertaken recently for environmental coastal monitoring using RS data and GIS \cite{BISHOPTAYLOR2021112734} The aim of these and similar works is to detect crisis areas where landscape changes take place, to spot and to map the affected regions for decision taking. For instance, human-induced processes may include urbanisation of coastal areas \cite{ANTOS2007173}, agriculture plantations \cite{6723525}, unsustainable tourism \cite{WESTON200998,Howesetal} or climate-related processes affecting ecosystems \cite{Glenn,Lemenkova202221}. These processes lead to the \hl{disruption of} natural ecosystems\hl{. They also may trigger changes in} land cover types \cite{DPCOSTA2024101574}, increase \hl{in} landscape fragmentation \cite{Lemenkova202406}, \hl{destruction of} connectivity between the individual patches and cause deforestation. In the end, such effects misbalance coastal ecosystem and modify its structure.  

Technically, detecting changes can be performed using time series analysis of the RS data\hl{. This is possible }through \hl{the} comparison of multi-temporal RS data covering the same area with different time intervals \cite{GAO2024119622,Vasquez2023,MASRIA2024103502}. Spatial analysis supported by\hl{ }multi-temporal RS data can be \hl{further} used to monitor and operatively detect such regions visible on the satellite images from space. \hl{Nevertheless, the} question arises\hl{ }how to effectively processes RS data and which methods \hl{among them} present the most efficient \hl{and optimised} solutions. The traditional tools of Geographic Information System (GIS) have restricted capabilities and limitations in speed and accuracy of RS data processing which is reported in relevant works \cite{BLACKFORD1995247,Whig2024,Rolim2023}. \hl{Therefore, }despite the proposed methodology\hl{, new technical approaches are} constantly developing (e.g., \cite{BUGNOT2018939}). 

Novel methods and approaches \hl{arise} in the rapidly developing \hl{era} of \hl{high-performance} computing\hl{. This enables} to overcome the \hl{existing} issues of RS data processing\hl{ and satellite image classification. For instance,} these include\hl{ }the use of \hl{the }general-purpose programming languages \hl{in cartographic applications, }such as Python or R, applications of Deep Learning (DL) \cite{Park,Roselin}, scripting techniques \cite{Lemenkova202212,Lemenkova202213} and algorithms of Artificial Intelligence (AI) \cite{rs11060617,Lemenkova202211}. The integration of RS data and these new technologies provides effective approach to RS data handling for accurate and rapid image processing \cite{Turnbull,Lemenkova202323}. \hl{Such} powerful combination of Earth observation data and ML tools\hl{ }presents a principally new cartographic approach for mapping marine and coastal environments\hl{ which can be effectively used }for \hl{operative }monitoring of marine and coastal areas \cite{Cao2022}.

Discriminating complex and highly fragmented patches in the coastal landscapes requires the use of the advanced methods for RS data processing. As a response to these needs, we utilise the algorithms of AI that enable to accurately recognise and classify even tiny differences in landscape patches and to detect the dynamics of vegetation types. Moreover, the workflow repeatability is facilitated using scripting techniques of \hl{Geographic Resources Analysis Support System (}GRASS) GIS and \hl{Generic Mapping Tools (}GMT\hl{)}. Such combined use of different software ensured the flexibility of the cartographic techniques through the use of GMT for gridded raster data handling using available techniques \cite{Lemenkova202217}, QGIS for vector data visualization and scripts of GRASS GIS for AI and ML methods of image processing. 

\hl{The} study is\hl{ }organised in the two \hl{logical }parts. \hl{The} first \hl{part} gives insights into the environmental setting of the coastal\hl{ }around the Phillip Port Bay, southern Australia, while the second one is \hl{focused on describing and presenting} technical the tools\hl{. Here we describe the details of the} methods of the AI algorithms \hl{which were }used for satellite image processing \hl{and point at existing differences} between \hl{them}. The results obtained from the image classification are\hl{ }commented in\hl{ relevant} sections with the sequential discussion, closing remarks and recommendations for future works. 

%----------- section ----------->
\section{Study Area}

The research is focused on the \hl{area} around Melbourne - the second largest city on the Australian continent, and encompasses the Port Phillip Bay. The exact zone of the study area is \hl{restricted between the coordinates} from 143°01'33.38" E to 143°08'03.62"E longitude and from 38°32'14.93"S to 36°24'10.94"S latitude\hl{. This area} is located in the southern part of the state of Victoria, \hl{Australia} \hl{(}Figure \ref{fig01}\hl{)}. 
 
 \begin{figure}[H]
    \begin{center}
	\hspace{-35pt}\includegraphics[width=14.0cm]{Fig_01.jpg}
    \end{center}
    \caption{General topographic map of Australia with outlined location of the study area. Mapping software: Generic Mapping Tools (GMT) version 6.4.0. Data source: GEBCO. Map source: author.}
    \label{fig01}
\end{figure}

The state of Victoria comprises \hl{of} 79 municipalities. \hl{According to Census 2023, Victoria had a population of 6,865,400} \cite{ABS}, which makes it the second-most-populated state (after New South Wales) in Australia. Such population density increases anthropogenic pressure on the natural \hl{environment} of coastal areas. On the other hand, there is a strong link between the coastal environment and quality of life for population living in this area which illustrates dense human-nature interactions \cite{LIAO2023111158,HUANG2020134680}. One of the most prominent features in the environment of southern Victoria is the coastal plains of Port Phillip Bay -- a largest coastal lagoon system in Victoria, Australia \hl{(}Figure \ref{fig02}\hl{)}. Port Phillipe Bay \hl{presents} \hl{an} embayment formed between the end of the last Ice Age around 8000 BCE as a result of faulting and marine transgression. Nowadays, it is located in a shallow tectonic depression with the depths mostly less than 8 m \cite{Holdgate}. With the total area of $1930 km^2$, it \hl{forms} a closed\hl{ }ecosystem with irregular coastline\hl{, diversified} topographic setting and \hl{complex }geomorphology \cite{longmore1996nutrient}. 

\begin{figure}[H]
    \begin{center}
	\hspace{-35pt}\includegraphics[width=15.0cm]{Fig_02.jpg}
    \end{center}
    \caption{Enlarged fragment of the topographic map of southern Australia showing the location of the study area. Mapping software: Generic Mapping Tools (GMT) version 6.4.0. Hydrographic and topographic data source: GEBCO. Map source: author.}
    \label{fig02}
\end{figure}

The dominating land cover types around \hl{the }Port Phillipe Bay and southern Victoria include bushland, native shrubland, pastures and grasslands that intersperse with native forests, Figure \ref{fig03}. Arid zones located further from the coastal area are notable for drylands with sparse crops and rare treed vegetation coverage with scattered trees. Lacustrine and moderate dune systems are also found in the surroundings of Port Phillip \hl{Bay}. The land affected by the anthropogenic activities includes agricultural farmlands, \hl{areas occupied by }horticulture and irrigated \hl{lands}, lands used for recreation, as well as built up areas. Those are \hl{mostly presented} by cities and \hl{urban }settlements, including small towns and\hl{ constructed artificial objects} (roads, industrial facilities, etc.) \hl{(}Figure \ref{fig03}\hl{)}.

\begin{figure}[H]
    \begin{center}
	\hspace{-35pt}\includegraphics[width=14.0cm]{Fig_03.jpg}
    \end{center}
    \caption{Land cover types in Australia. Mapping software: QGIS version 3.34.7 'Prizren'. Data source: Dynamic Land Cover Dataset (DLCD). Map source: author.}
    \label{fig03}
\end{figure}

\hl{Landscape formation} in this area is driven by diverse factors such as \hl{geologic foundations, climate and hydrological processes, e.g., }wind force,\hl{ }strong effects of oceanic waves, tides, and currents. The variability of such processes affects nearby landscapes, controls the distribution of the land cover types and regulates \hl{general }ecosystem functioning. As a result, the landscapes surrounding \hl{ the estuaries of the }Port Phillip Bay are presented by a complex mixture of \hl{types, including} wetlands\hl{. Wetlands }create a dynamic environment which supports rich marine and coastal biodiversity and important fisheries \cite{Jenkins,Weston,harris1996port}.

\hl{Recent} environmental survey \hl{in the southern Victoria around the Port Phillip Bay }reported \hl{the distribution of }19212 ha of coastal salt marshes 5177 ha of mangroves, and 3227 ha of estuarine wetland in \hl{this area} \cite{Boon}. Among these, the most notable system of wetlands is Cheetham, which occupies up to 420 ha of lagoons, including both artificial (salt marshes) and natural types (mangroves) \hl{(}Figure \ref{fig02}\hl{)}. \hl{High urbanization rate and constantly increasing population density} of the Port Phillip Bay \hl{is especially notable along the southern} coasts of Victoria\hl{, }within 100 km of coastlines\hl{. As a result, anthropogenic} pressure, exaggerated by \hl{industrial} activities along the coasts, \hl{lead to the} degradation of coastal areas and deterioration of the landscapes around wetland areas \cite{land11020311}\hl{. Other environmental consequences include declined} water quality, \hl{exceeded} runoff of agricultural fertilizers, pesticides and chemical into the Port Philippe Bay. \hl{Chemical pollution }affected water quality \hl{through} high suspended sediments with nutrient inflows \hl{which caused} eutrophication. 

Cheetham Wetlands are distributed on old salt works land on the western part of Port Phillip Bay and include the seasonal and perennial \hl{types} protected for conservation purposes. Located ca. 20 kilometres southwest of Melbourne \cite{Westonetal2006}, perennial Cheetham Wetlands dry occasionally, only during periods with low precipitation \cite{jmse9080898}. However, much of the flora and fauna are not adapted to drying and remain permanently in the areas filled with water. Such settings create unique environmental ecosystem\hl{ }with specific species adapted to such conditions of variable water levels. Hydrological and climate setting of the Cheetham Wetlands favour its high environmental value by providing habitat over 200 species of rare birds\hl{. Apart from regional species, they also include} migratory \hl{birds} that \hl{move} to the Southern Hemisphere between July and November \cite{Jackson}. 

The current environment of Port Phillip Bay has developed by a complex interactions of climate, lithological variability, tectonic-geologic deformation and sea level changes. Nowadays, \hl{regional} land cover types\hl{ }continue further changing both along the coastline and in the hinterland\hl{. This} is mainly caused by human activities such as economic development, residential construction\hl{, }urban sprawl \hl{and} recreational pressures \cite{FultonSmith}. Current environmental issues in this area include anthropogenic threats such as commercial dredging, water pollution and contamination caused by the trace metals, toxic substances, organochlorines and hydrocarbons \cite{PHILLIPS1992200}. \hl{Recent} studies also detected changes \hl{in} structure and composition of benthic communities \hl{as well as} weed invasions in Port Phillip Bay caused by the environmental interactions \cite{CURRIE199936}. Finally, shallow bathymetry of the \hl{bay} contributes to the eutrophication \cite{FultonSmith} and marine algae bloom \cite{md18030142,976565} in the selected parts of the embayment. As a result, the overall condition of Port Phillip Bay appears to have deteriorated on a large scale over the recent decades \cite{JUNG201193}.

%----------- section ----------->
\section{Materials and Methods}

An approach that integrates RS data and AI algorithms is proposed herein to identify land cover changes through pixel-level mapping for \hl{environmental }monitoring\hl{ of the }southern area of Victoria state, Melbourne district. 

%----------- subsection ----------->
\subsection{Data}

The data used in this study for spatial analysis and image processing \hl{were} obtained from the open sources. The geospatial data of GEBCO (\href{https://www.gebco.net/}{https://www.gebco.net/}, accessed on 24 June 2024) with a 15 arc second spatial resolution was incorporated for topographic mapping. The general topographic map with the outlined location of the study area within Australia was plotted based on the General Bathymetric Chart of the Oceans (GEBCO) dataset with 15 arc minute resolution, as shown in Figures \ref{fig01} and \ref{fig02}. The names of the major features, cities and geographic objects were obtained from the Gazetteer of Australia (Aggregator: Geoscience Australia, distributed under the Creative Commons Attribution 3.0 license). 

The\hl{ }data\hl{ for the land cover types reference }were obtained from the open repository of the Geoscience Australia of the Australian Government\hl{. In this dataset, these categories are identified }based on the updated classified imagery of Geoscience Australia Landsat Land Cover with resolution of 25\hl{ }m that replaces the previous version based on classified images from Terra Moderate Resolution Imaging Spectroradiometer (MODIS). \hl{Nowadays, it presents }the nationally consistent and thematically comprehensive land cover reference for Australia -- the Dynamic Land Cover Dataset (DLCD)\hl{, }version 1.0.0 (\href{https://knowledge.dea.ga.gov.au/data/product/dea-land-cover-landsat/?tab=overview}{https://knowledge.dea.ga.gov.au/data/product/dea-land-cover-landsat/?tab=overview}, accessed on 24 June 2024), Figure \ref{fig03}. These data were used to obtain the general characteristics of the landscape types in the Port Phillip Bay and the surroundings\hl{. }The vectors layers of these data were processed and visualised using QGIS software.\hl{ }

\hl{Coping with the diversity of information in the RS data with different spatio-temporal or categorical attributes poses a prevalent difficulty in satellite images processing. With this regards, selecting data with suitable spatio-temporal parameters play a crucial role in mitigating this challenge by furnishing tailored suggestions. Thus, environmental analysis takes advantage of diverse aspects of RS such as physical fundamentals of spectral reflectance, orbit characteristics that differ various satellite missions. To compare some of them, a single scene of Satellite Pour l'Observation de la Terre (SPOT) image covers a footprint of} $3,600 km^2$ \hl{at resolutions of 20 m to 2.5 m, with a location accuracy up to 10 m, Landsat 8-9 OLI/TIRS imagery have resolution of 30 m in multispectral bands, Sentinel-2 imagery has different resolution in four bands at 10 m, six bands at 20 m and three bands at 60 m spatial resolution. }

\hl{Spatial data variability and difference of their technical characteristics illustrate by the examples above require adjusted approaches, such as ML and AI, to handle these data in order to obtain geo-information effectively. Therefore, }the RS data \hl{used in this study }were obtained from the sources of Landsat Operational Land Imager and Thermal Infrared Sensor (OLI-TIRS) United States Geological Survey (USGS)\hl{. The data were selected due to their high quality, acceptable resolution, robust technical characteristics and cloudless coverage of the study area}. The main metadata of the satellite images are summarised in Table \ref{tab01}. 

\begin{sidewaystable}
\centering
\small{
 \begin{tabular}{|l|l|l|l|l|}
    \hline
        \textbf{Data Set Attribute} & \textbf{Attribute Value: 2024} & \textbf{Attribute Value: 2017} & \textbf{Attribute Value: 2015} & \textbf{Attribute Value: 2013} \\ \hline
        \toprule
        Landsat Product Identifier L2 & \makecell{LC09\_L2SP\_093086\_ \\ 20240329\_20240403\_02\_T1} & \makecell{LC08\_L2SP\_093086\_ \\ 20170318\_20200904\_02\_T1} & \makecell{LC08\_L2SP\_093086\_ \\ 20150329\_20200909\_02\_T1} & \makecell{LC08\_L2SP\_093086\_ \\ 20130326\_20200913\_02\_T1} \\ \hline
        Landsat Product Identifier L1 & \makecell{LC09\_L1TP\_093086\_ \\ 20240329\_20240329\_02\_T1} & \makecell{LC08\_L1TP\_093086\_ \\ 20170318\_20200904\_02\_T1} & \makecell{LC08\_L1TP\_093086\_ \\ 20150329\_20200909\_02\_T1} & \makecell{LC08\_L1TP\_093086\_ \\ 20130326\_20200913\_02\_T1} \\ \hline
        Landsat Scene Identifier & LC90930862024089LGN00 & LC80930862017077LGN00 & LC80930862015088LGN01 & LC80930862013085LGN02 \\ \hline
        Date Acquired & 2024/03/29 & 2017/03/18 & 2015/03/29 & 2013/03/26 \\ \hline
        Roll Angle & 0.001 & 0.000 & 0.000 & 0.000 \\ \hline
        Date Product Generated L2 & 2024/04/03 & 2020/09/04 & 2020/09/09 & 2020/09/13 \\ \hline
        Date Product Generated L1 & 2024/03/29 & 2020/09/04 & 2020/09/09 & 2020/09/13 \\ \hline
        Start Time & 2024-03-29 00:09:22 & 2017-03-18 00:09:03.868153 & 2015-03-29 00:08:50.482423 & 2013-03-26 00:10:37.681759 \\ \hline
        Stop Time & 2024-03-29 00:09:54 & 2017-03-18 00:09:35.63815 & 2015-03-29 00:09:22.252419 & 2013-03-26 00:11:07.47778 \\ \hline
        Land Cloud Cover & 0.34 & 0.26 & 0.48 & 0.13 \\ \hline
        Scene Cloud Cover L1 & 0.33 & 0.24 & 0.45 & 0.12 \\ \hline
        Ground Control Points Model & 1160 & 1194 & 1171 & 1081 \\ \hline
        Geometric RMSE Model & 5.002 & 4.048 & 4.001 & 4.701 \\ \hline
        Geometric RMSE Model X & 3.304 & 2.534 & 2.419 & 3.135 \\ \hline
        Geometric RMSE Model Y & 3.756 & 3.156 & 3.188 & 3.503 \\ \hline
        Processing Software Version & LPGS\_16.4.0 & LPGS\_15.3.1c & LPGS\_15.3.1c & LPGS\_15.3.1c \\ \hline
        Sun Elevation L0RA & 38.05564986 & 41.13467347 & 38.22640744 & 39.13419122 \\ \hline
        Sun Azimuth L0RA & 45.77175870 & 50.09114471 & 46.16829406 & 46.70481697 \\ \hline
        TIRS SSM Model & N/A & FINAL & ACTUAL & ACTUAL \\ \hline
        Scene Center Latitude & -37.47464 & -37.47435 & -37.47440 & -37.47462 \\ \hline
        Scene Center Longitude & 144.35289 & 144.40120 & 144.39626 & 144.31862 \\ \hline
        Corner Upper Left Latitude & -36.40217 & -36.40099 & -36.40088 & -36.45905 \\ \hline
        Corner Upper Left Longitude & 143.10763 & 143.15450 & 143.15116 & 143.11148 \\ \hline
        Corner Upper Right Latitude & -36.45824 & -36.45609 & -36.45602 & -36.51418 \\ \hline
        Corner Upper Right Longitude & 145.66861 & 145.71886 & 145.71217 & 145.59395 \\ \hline
        Corner Lower Left Latitude & -38.47670 & -38.48103 & -38.48092 & -38.42299 \\ \hline
        Corner Lower Left Longitude & 142.99849 & 143.04638 & 143.04295 & 143.00833 \\ \hline
        Corner Lower Right Latitude & -38.53712 & -38.54041 & -38.54033 & -38.48215 \\ \hline
        Corner Lower Right Longitude & 145.63117 & 145.68274 & 145.67586 & 145.55655 \\ \hline
    \end{tabular}
    }
\caption{Metadata for the four scenes of the Landsat 8-9 OLI/TIRS images used for RS data processing and classification}
\label{tab01}
\end{sidewaystable}

Technically, the multispectral Landsat 8-9 OLI-TIRS images were retrieved from the Earth Explorer repository (\href{https://earthexplorer.usgs.gov/}{https://earthexplorer.usgs.gov/}, accessed on 24 June 2024). The metadata were acquired and checked to cloudiness coverage and image quality and image extent to cover the entire study area of the Port Phillip Bay. Four satellite images were obtained for the following dates in \hl{the month of }March: 26/03/2013, 29/03/2015, 18/03/2017 and 29/03/2024. The pictures of these images as a raw dataset in natural colours are presented in Figure \ref{fig04}.

\begin{figure}[H]
  \centering
  \begin{tabular}[b]{c}
    \includegraphics[width=2.9cm]{Fig_04a.jpg} \\
    \small (a): March 2013
  \end{tabular}
  \begin{tabular}[b]{c}
    \includegraphics[width=2.9cm]{Fig_04b.jpg} \\
    \small (b): March 2015
  \end{tabular}
    \begin{tabular}[b]{c}
    \includegraphics[width=2.9cm]{Fig_04c.jpg} \\
    \small (c): March 2017
  \end{tabular}
  \begin{tabular}[b]{c}
    \includegraphics[width=2.9cm]{Fig_04d.jpg} \\
    \small (d): March 2024
  \end{tabular}
  \caption{Original data: Landsat 8-9 OLI/TIRS images on southern Australia, Philippe Bay collected during March on years 2013, 2015, 2017 and 2024. Source: USGS. Compilation source: author.}
  \label{fig04}
\end{figure}

Major technical characteristics common for all the Landsat scenes are the following\hl{: the} satellite images have a level Data Type L2, meaning they incorporate geometric correction and topographic correction based on parameters of OLI TIRS L2SP sensor identifier, collection category T1, collection number 2 and Ground Control Points (GCP) Version 5. Special resolution is 30 m for multispectral channels and cirrus (Bands 1 to 7 and Band 9), 15 m for Panchromatic (Band 8) and 90 m for Thermal Infrared Sensor (TIRS) (Bands 10 and 11). The cartographic parameters are the following\hl{: the} data are projected in Universal Transverse Mercator (UTM), Zone 55 South, Datum and Ellipsoid WGS84. The Satellite ID is 9 for the scene on 2024 and 8 for the rest of the images, and the Station Identifier is LGN for all the data. All images were taken in nadir in day time. The Worldwide Reference System (WRS) Path is 93, Row 86. 

%----------- subsection ----------->
\subsection{Workflow}

The workflow steps to process and classify the satellite images in order to detect land cover types along the coastal \hl{landscapes of} southern Australia are \hl{schematically presented} in Figure \ref{fig05}. 

\begin{figure}[H]
    \begin{center}
	\hspace{-35pt}\includegraphics[width=13.0cm]{Fig_05.jpg}
    \end{center}
    \caption{Workflow methodology. Software: RStudio version 2024.04.2+764, R version 3.6.0, 'DiagrammeR' package version 1.0.11.9000. Diagram source: author.}
    \label{fig05}
\end{figure}

\hl{Employing ML and AI techniques of GRASS GIS involves a multi-step process. The first stage involves the data import through the modules 'r.import' or 'r.in.gdal'. The formulation of the initial data check includes metadata control using 'r.info' and 'r.category' modules which provide information on raster map layer and evaluates the values and labels of categories, respectively. Among others, it enables to check the background information about the geographic extent of the study area and variables in an analysis of raster layer. The g.copy module copies the existing raster maps in the mapset which is controlled by the module 'g.list rast' for evaluating its current content.}

\hl{During the second stage (image processing), GRASS GIS begins to explore the images by generating corrected scenes. To this end, the atmospheric correction, geometric rectification and calibration are applied using the modules 'i.landsat.toar', 'i.rectify'  and others.}\hl{At this stage, the quality of the multispectral bands of Landsat is improved to facilitate the next steps of image classification by the AI and ML algorithms. Cartographic handling enables to tune up the region extent, visualise the data, add explanatory legend, create colore composites from the multispectral triplets, and select representative colours for each map using several modules: 'd.rast', 'd.legend', and 'r.colors'. In the next stage, GRASS GIS narrows the focus of image processing to the ML algorithms (RF, SVM, MLP, etc.) and evaluates the values cells using computer vision approach. The pixels are automatically distributed and grouped into classes according to their DNs.} 

\hl{This process involves the use of modules 'r.random', 'r.learn.train' and other algorithms, as presented in the diagram scheme (Figure} \ref{fig05}\hl{). The ML part is performed using algorithms of Scikit-Learn package of Python that is integrated with GRASS GIS, version 8.3. In the final stage, the classification of the images is finalised with map outputs and accuracy assessment. The latter includes the computed rejection probability accuracy using chi-square algorithm, generating statistical report, and computing class separability matrices for each year of the evaluated Landsat scenes. At this point, the quality of processed data is assed using particular algorithms embedded into GRASS GIS.}.

The adjusted automatic classification model was performed using the "i.maxlik" module of GRASS GIS that employs the maximum-likelihood discriminant analysis classifier. \hl{The thematic map on land cover types was created using QGIS software to evaluate the distribution of major categories across the study area. Additionally, a digital elevation model (DEM) was obtained from GEBCO and generated as a topographic map of study area to strengthen the environmental setting through geospatial information. }Here, the initial values of the pixels were computed for each multispectral band of the Landsat scenes for each year (the results are reported in the Appendix \ref{AppA}). Then, the iterative process of the re-assignment of these pixels was performed (the results of the iteration cycles are reported in the Appendix \ref{AppB} for each year), and the pixels were harmonised and readjusted to the target classes through the repetitive sequence of computational process. 

The cluster parameters for MaxLike classification approach were accepted \hl{with the }following \hl{values}. The number of initial classes is 10; the minimum class size is 17; the Minimum class separation is zero and the percent convergence is 98.0 for all the scenes. The maximum number of iterations is 30 as defined by the GRASS GIS algorithm. The statistical summary of the accepted computational parameters of image processing through clustering is given in Table \ref{tab02}.

\begin{table}[H]
  \centering
  \begin{tabularx}{\textwidth}{p{0.9cm}p{0.9cm}p{0.9cm}p{1.0cm}p{2.0cm}p{2.2cm}p{2.5cm}}
  \toprule
  \multirow{2}{*}{Year} & \multicolumn{6}{c}{Parameters of clustering} \\
     & Rows & Cols & Cells & Sample size & Row sampling interval & Column sampling interval \\
    \midrule
    \textbf{2013} & 7281 & 7421 & 54032301 & 7044 & 72 & 74 \\
    \textbf{2015} & 7711 & 7661 & 59073971 & 7005 & 77 & 76 \\
    \textbf{2017} & 7711 & 7652 & 59004572 & 7011 & 77 & 76 \\
    \textbf{2024} & 7692 & 7522 & 57859224 & 7217 & 76 & 75 \\
\bottomrule
\end{tabularx}
  \caption{Statistical parameters for clustering using 'i.maxlik' module of GRASS GIS}
  \label{tab02}
\end{table}

\hl{Four widely used ML algorithms were applied to classify the images: 1) Random Forest (RF), 2) Support 7 Vector Machine (SVM), 3) ANN-based approach using Multi-layer Perceptron (MLP) Classifier; and 4) Decision Tree Classifier (DTC). The SVM is a non-parametric statistical learning classifier which is initially proposed by Vapnik} \cite{788640} \hl{with the aim to find the optimal surface that divides each set of points into distinct classes according to their properties. The optimal surface here means a decision boundary that minimises the number of erroneously classified points. Such logical approach increases the accuracy of the classification and which resulted in the high applicability of this method in image classification and explains its high accuracy.}

\hl{The} maps obtained based on the classified images were visualised using the sequence of modules "d.rast", "d.grid" and "d.legend" \hl{(}Figure \ref{fig06}\hl{)}.  Such workflow made it possible to automate the necessary procedures to carry out image analysis. The additional information from the spectral bands (one to seven)  was obtained and statistically processed to evaluate changes during the period from 2013 to 2024. 

Since all the four ML and \hl{AI} models that were applied belong to the type of supervised classifiers, a training dataset was necessary for the algorithm. Training dataset teaches the model to learn the patterns of the different classes which should be discriminated. To this end, a sample of the unsupervised classification with pixel points was collected from the clustering to train the ML and AI algorithms, Figure \ref{fig07}. This initial raster dataset was augmented through the visual recognition of the satellite images, using the data on land cover types \hl{which were }collected as reference from the \hl{Dynamic Land Cover Dataset (}DLCD) open repository of the Australian Government. 

\hl{Four} different ML and AI algorithms were tested and their performances \hl{were} evaluated. For SVM and RF, the optimal parameters of the models were adjusted to improve the classification performance and accuracy of assignment. At the same time, for all the years (2013-2024), the assignment of pixels to various land cover classes was evaluated for each band was calculated to understand the contribution of land cover types within the whole landscape in the classification models.  

%----------- section ----------->
\section{Results}

\subsection{Classification output}

\hl{The results} of the classification of the satellite images using ML and AI methods are presented in Figures \ref{fig06}, \ref{fig07}, \ref{fig08} and \ref{fig09} for comparison of different approaches for years 2013, 2015, 2017 and 2024. Contrariwise, the time series of one algorithms but evaluating the environmental dynamics is illustrated in Figure \ref{fig10}. The computed statistics on land cover class distribution by multispectral bands for each Landsat image and by years for each evaluated period (2013, 2015, 2017 and 2024) \hl{were} obtained and summarised in the tables presented in the following section as well as in the Appendix of this study, e.g., Table \ref{tab03} for 2013 and likewise for the next periods. For 2013, pixels were assigned to the corresponding categories of the recognised 10 classes with 98.27\% points stable accuracy and convergence of 98.3\%, Table \ref{tab03}.

\hl{Figure} \ref{fig07} \hl{and Figure} \ref{fig08} \hl{displays the results of the image classification approach using four different AI and ML methods by GRASS GIS resulting in ten land cover classes as described in using seven multispectral bands from the 2015 and 2017 Landsat images. Visually, these maps reflect broad LandUse Land Cover (LULC) characteristics of the Cheetham Wetlands, section of the marsh and uplands surrounding Port Phillip Bay.} \hl{The Cheetham Wetlands and marsh area (Class 7) are bounded to the west by upland vegetation including mixed tree cover and sparse vegetation and scattered trees (Class 5), as well as scrubs, bushland and shrubland (Class 2). The marsh areas and wetlands are partially mingled with other cover classes and inclusions of sigle trees and bushes adopted to the wet conditions. }

\begin{figure}[H]
  \centering
  \begin{tabular}[b]{c}
    \includegraphics[width=6.4cm]{Fig_06a.jpg} \\
    \small (a)
  \end{tabular}
  \begin{tabular}[b]{c}
    \includegraphics[width=6.4cm]{Fig_06b.jpg} \\
    \small (b) 
  \end{tabular}
    \begin{tabular}[b]{c}
    \includegraphics[width=6.4cm]{Fig_06c.jpg} \\
    \small (c)
  \end{tabular}
  \begin{tabular}[b]{c}
    \includegraphics[width=6.3cm]{Fig_06d.jpg} \\
    \small (d) 
  \end{tabular}
  \caption{Results of image processing on 2013 using four methods: 1) Random Forest; 2) SVM; 3) DecisionTreeClassifier; 4) ANN-based MLPClassifier. Image processing: author.}
  \label{fig06}
\end{figure}

\hl{The} selection of the\hl{ }land categories has been employed using adopted \hl{data from} southern Australia and Port Phillip Bay. The adaptation of these data aimed at the correct detection of coastal landscapes within the ecosystems of southern Victoria with the special aim to detect wetlands, and discriminate them automatically from urban areas and native forests. \hl{The }retrospective time series 1987-2019 over Victoria State \hl{was employed }(\href{https://www.environment.vic.gov.au/biodiversity/Victorias-Land-Cover-Time-Series}{https://www.environment.vic.gov.au/biodiversity}, accessed on 25 July 2024): \emph{1) water bodies (Port Phillip Bay, lakes, rivers); 2) bushland ; 3) farmland and cropland used for agriculture; 4) native forests (closed trees); 5) sparse vegetation and scattered trees; 6) recreational land; 7) wetlands; 8) built-up areas; 9) pastures and grasslands; 10) bare areas}\hl{.} 

\hl{To the east along the Indian Ocean, developed land appears as a segment immediately adjacent to the open coastal spaces. This includes a section of marshes beachfront development notable by real estate of compact beach front constructions close to the water border. The classes of natural vegetation are also observed in other segments of the classified scene and includes features such as native forests (closed trees) and sparse vegetation and scattered trees in drylands. The regions marked by anthropogenic activities (Class 8) include built-up areas which comprise the constructed commercial and residential buildings, roads and adjacent parking places and related objects of infrastructure. The distribution of the farmland and cropland (Class 3) is notable in the eastern areas with more intense agriculture activities.}

\hl{The shallow water in the Port Phillip Bay are distinct in colour from the waters of the Indian and Pacific Oceans, and are typified by very low lying areas near river and creek edges, and also marked by existing pools, wet soils along the coasts, and low marsh vegetation along creek edges that belong to the hydrological system of Cheetham Wetlands. Likewise, the class of coastal marshes and wetlands also contain a few smaller pools and marginally wet areas where specific types of vegetation develop.} 

\begin{figure}[H]
  \centering
  \begin{tabular}[b]{c}
    \includegraphics[width=6.4cm]{Fig_07a.jpg} \\
    \small (a)
  \end{tabular}
  \begin{tabular}[b]{c}
    \includegraphics[width=6.4cm]{Fig_07b.jpg} \\
    \small (b) 
  \end{tabular}
    \begin{tabular}[b]{c}
    \includegraphics[width=6.4cm]{Fig_07c.jpg} \\
    \small (c)
  \end{tabular}
  \begin{tabular}[b]{c}
    \includegraphics[width=6.4cm]{Fig_07d.jpg} \\
    \small (d) 
  \end{tabular}
   \caption{Image processing on 2015 using four methods: 1) Random Forest; 2) SVM; 3) ANN-based MLPC; 4) DTC. Image source: author.}
  \label{fig07}
\end{figure}

\begin{table}[H]
  \centering
  \tiny
  \begin{tabularx}{\linewidth}{p{1.3cm}p{1.0cm}p{1.0cm}p{1.0cm}p{1.0cm}p{1.0cm}p{1.0cm}p{1.0cm}p{1.0cm}}
  \toprule
  \multirow{2}{*}{Class} & \multicolumn{7}{c}{Multispectral bands of the Landsat 8 OLI/TIRS scene for \textbf{2013}} \\
    & Stat & B1 & B2 & B3 & B4 & B5 & B6 & B7 \\
    \midrule
    \multirow{2}{*}{1 (507)} & means & 7687.78 & 7908.5 & 8088.38 & 7475.44 & 7336.59 & 7412.26 & 7391.69 \\
    & stdev & 344.653 & 421.512 & 668.419 & 537.547 & 460.436 & 295.948 & 180.949 \\
    \midrule
    \multirow{2}{*}{2 (638)} & means & 7680.67 & 7842.55 & 8326.59 & 8342.63 & 14223.1 & 10929.8 & 9067.1 \\
    & stdev & 222.242 & 267.302 & 375.086 & 444.833 & 1553.37 & 856.04 & 567.558 \\
    \midrule
    \multirow{2}{*}{3 (887)} & means & 8097.13 & 8335.14 & 8934.32 & 9236.08 & 14634.2 & 13461.6 & 10833.1 \\
    & stdev & 256.84 & 273.675 & 349.578 & 379.762 & 1150.47 & 761.219 & 550.821 \\
    \midrule
    \multirow{2}{*}{4 (617)} & means & 8656.62 & 9001.7 & 9802.38 & 10300 & 14321.4 & 14803.1 & 12391.7 \\
    & stdev & 610.8 & 666.747 & 803.156 & 815.317 & 969.65 & 962.333 & 703.455 \\
    \midrule
    \multirow{2}{*}{5 (378)} & means & 8457.53 & 8775.73 & 9744.8 & 10004.9 & 17329.6 & 16009.9 & 12429.8 \\
    & stdev & 635.538 & 594.932 & 704.354 & 866.12 & 1584.96 & 1076.47 & 793.116 \\
    \midrule
    \multirow{2}{*}{6 (748)} & means & 8753.46 & 9158.45 & 10033.5 & 10785.5 & 14633.9 & 17307.8 & 13990.3 \\
    & \textbf{stdev} & 306.286 & 344.408 & 426.607 & 563.988 & 788.463 & 802.195 & 664.108 \\
    \midrule
    \multirow{2}{*}{7 (584)} & means & 8868.88 & 9316.54 & 10365.6 & 11059.3 & 16695 & 18546 & 14492.1 \\
    & stdev & 240.485 & 270.356 & 368.04 & 570.249 & 1012.67 & 773.668 & 561.7 \\
    \midrule
    \multirow{2}{*}{8 (1004)} & means & 9161.49 & 9674.29 & 10660.8 & 11716.9 & 15122.4 & 19256.5 & 15629.3 \\
    & stdev & 354.297 & 391.741 & 494.712 & 625.917 & 665.09 & 805.656 & 665.977 \\
    \midrule
    \multirow{2}{*}{9 (1175)} & means & 9392.39 & 9999.52 & 11149.5 & 12418.8 & 16298.3 & 20847.9 & 16690.6 \\
    & stdev & 251.846 & 295.699 & 411.786 & 635.913 & 797.684 & 750.01 & 752.185 \\
    \midrule
  \multirow{2}{*}{10 (506)} & means & 9967.75 & 10719.8 & 12141.4 & 13725.9 & 17504.8 & 22182.4 & 18140.5 \\
   & stdev & 919.948 & 1000.56 & 1198.41 & 1274.36 & 1105.95 & 1361.48 & 1327.76 \\
\bottomrule
\end{tabularx}
  \caption{Pixelwise computational statistics on land cover class distribution by spectral bands on 2013.}
  \label{tab03}
\end{table}

\subsection{Statistical metrics}
Table \ref{tab04} shows the evaluated statistical metrics derived from the computations obtained for each multispectral class for 2015.

\begin{table}[H]
  \centering
  \begin{tabularx}{\linewidth}{p{1.3cm}p{1.0cm}p{1.0cm}p{1.0cm}p{1.0cm}p{1.0cm}p{1.0cm}p{1.0cm}p{1.0cm}}
  \toprule
  \multirow{2}{*}{Class} & \multicolumn{7}{c}{Multispectral bands of the Landsat 8 OLI/TIRS scene for \textbf{2015}} \\
    & Stat & B1 & B2 & B3 & B4 & B5 & B6 & B7 \\
    \midrule
    \multirow{2}{*}{1 (568)} & means & 7940.28 & 8075.57 & 8123.99 & 7524.04 & 7425.85 & 7399.9 &7372.91 \\
    & stdev & 344.325 & 379.953 & 593.974 & 522.724 & 652.721 & 335.877 & 246.953 \\
    \midrule
    \multirow{2}{*}{2 (615)} & means & 7686.57 & 7824.36 & 8272.39 & 8270.83 & 13898.8 & 10527 & 8863.34 \\
    & stdev & 305.668 & 359.903 & 482.018 & 521.82 & 1537.46 & 845.958 & 552.008 \\
    \midrule
    \multirow{2}{*}{3 (781)} & means & 8072.26 & 8269.64 & 8837.17 & 9086.89 & 14572.4 & 13012 & 10576.3 \\
    & stdev & 259.275 & 259.011 & 320.839 & 364.833 & 1119.29 & 746.203 & 547.4 \\
    \midrule
    \multirow{2}{*}{4 (703)} & means & 8650.11 &8941.9 & 9661.26 & 10136 & 14137.3 & 14558.4 & 12268.6 \\
    & stdev & 617.366 & 675.844 & 814.999 & 863.171 & 1032.01 & 978.397 & 738.674 \\
    \midrule
    \multirow{2}{*}{5 (391)} & means & 8462.31 & 8728.8 & 9605.56 & 9842.91 & 17552 & 15561.5 & 12141.8 \\
    & stdev & 293.506 & 297.3 & 383.437 & 489.254 & 1738.29 & 1098.23 & 795.876 \\
    \midrule
    \multirow{2}{*}{6 (796)} & means & 8861.33 & 9180.17 & 9941.97 & 10635.8 & 14647.6 & 17004.9 & 13981.4 \\
    & stdev & 356.821 & 369.241 & 445.523 & 545.923 & 900.559 & 805.354 & 717.999 \\
    \midrule
    \multirow{2}{*}{7 (582)} & means & 8969.33 & 9327.61 & 10291.4 & 10917.5 & 17000.9 & 18508.4 & 14546.8 \\
    & stdev & 476.562 & 385.029 & 458.457 & 582.109 & 1075.41 & 907.064 & 678.588\\
    \midrule
    \multirow{2}{*}{8 (1026)} & means & 9211.47 & 9575.03 & 10376.8 & 11287.4 & 14865.5 & 19153.9 & 15775.7 \\
    & stdev & 278.492 & 282.228 & 361.572 & 493.88 & 732.319 & 843.458 & 637.941 \\
    \midrule
    \multirow{2}{*}{9 (1245)} & means & 9452.33 & 9909.68 & 10884.4 & 11989 & 15914.4 & 21031.7 & 17155.4 \\
    & stdev & 237.238 & 250.904 & 380.325 & 578.133 & 859.887 & 805.656 & 762.815 \\
    \midrule
    \multirow{2}{*}{10 (298)} & means & 10110.3 & 10774.9 & 12172.6 & 13679.5 & 17597.3 & 22695 & 19191.3 \\
    & stdev & 1645.09 & 1738.58 &1908.73 & 1940.05 & 2142.28 & 1510.8 & 1592.87 \\
    \midrule 
\bottomrule
\end{tabularx}
  \caption{Pixelwise computational statistics on land cover class distribution by spectral bands of the Landsat-8 OLI/TIRS image on 2015}
  \label{tab04}
\end{table}

\begin{figure}[H]
  \centering
  \begin{tabular}[b]{c}
    \includegraphics[width=6.4cm]{Fig_08a.jpg} \\
    \small (a)
  \end{tabular}
  \begin{tabular}[b]{c}
    \includegraphics[width=6.4cm]{Fig_08b.jpg} \\
    \small (b) 
  \end{tabular}
    \begin{tabular}[b]{c}
    \includegraphics[width=6.4cm]{Fig_08c.jpg} \\
    \small (c)
  \end{tabular}
  \begin{tabular}[b]{c}
    \includegraphics[width=6.4cm]{Fig_08d.jpg} \\
    \small (d) 
  \end{tabular}
  \caption{Results of image processing on 2017 using four different methods: 1) Random Forest; 2) Support Vector Machine (SVM); 3) ANN-based approach using MLPClassifier; 4) DecisionTreeClassifier. Image processing source: author.}
  \label{fig08}
\end{figure}

Likewise, the performance of the image classification for years 2017 and 2024 were summarised in Table \ref{tab05} and Table \ref{tab06}. It shows the variations in pixels's means and standard deviation for pixels by Landsat channel. 

\begin{table}[H]
  \centering
  \begin{tabularx}{\linewidth}{p{1.3cm}p{1.0cm}p{1.0cm}p{1.0cm}p{1.0cm}p{1.0cm}p{1.0cm}p{1.0cm}p{1.0cm}}
  \toprule
  \multirow{2}{*}{Class} & \multicolumn{7}{c}{Multispectral bands of the Landsat 8 OLI/TIRS scene for \textbf{2017}} \\
    & Stat & B1 & B2 & B3 & B4 & B5 & B6 & B7 \\
    \midrule
    \multirow{2}{*}{1 (581)} & means & 7929.46 & 8108.94 & 8337.41 & 7709.38 & 7490.57 & 7424.3 & 7396.55 \\
    & stdev & 337.158 & 410.139 & 735.841 & 802.125 & 750.097 & 360.916 & 220.256 \\
    \midrule
    \multirow{2}{*}{2 (673)} & means & 7679.48 & 7836.62 & 8328.88 & 8317.19 & 14848.9 & 10861.6 & 8990.81 \\
    & stdev & 223.056 & 254.243 & 345.482 & 464.513 & 1369.21 & 883.749 & 560.842 \\
    \midrule
    \multirow{2}{*}{3 (494)} & means & 8537.22 & 8839.63 & 9573.69 & 9971.95 & 13978 & 13613.1 & 11552 \\
    & stdev & 633.576 & 671.098 & 776.416 & 867.675 & 1093.82 & 1071.02 & 823.154 \\
    \midrule
    \multirow{2}{*}{4 (674)} & means & 7962.61 & 8209.86 & 8883.57 & 9075.13 & 16186 & 13121.2 & 10398.7 \\
    & stdev & 199.388 & 190.559 & 260.072 & 330.17 & 1246.94 & 716.604 & 522.388 \\
    \midrule
    \multirow{2}{*}{5 (534)} & means & 8467.7 & 8822.3 & 9723.03 & 10205.7 &16661.4 & 15621.8 & 12327 \\
    & stdev & 459.728 & 463.622 & 524.131 & 604.138 & 1231.37 & 854.132 & 622.67 \\
    \midrule
    \multirow{2}{*}{6 (836)} & means & 8762.76 & 9155.02 & 10011 & 10828.8 & 14450.9 & 16939 & 13601.9 \\
    & stdev & 394.902 & 369.181 & 436.331 & 516.127 & 772.106 & 926.651 & 639.347 \\
    \midrule
    \multirow{2}{*}{7 (629)} & means & 8936.39 & 9422.97 & 10520.8 & 11467.8 & 16778.1 & 18160 & 14163.3 \\
    & stdev & 427.325 & 472.756 & 601.364 & 721.644 & 1091.77 & 848.855 & 622.946 \\
    \midrule
    \multirow{2}{*}{8 (1048)} & means & 9073.67 & 9536.34 & 10454.4 & 11506.1 & 14986.6 & 19085.9 & 15192 \\
    & stdev & 296.637 & 272.932 & 327.833 & 442.594 & 650.591 & 671.502 & 648.356 \\
    \midrule
    \multirow{2}{*}{9 (1140)} & means & 9304.19 & 9868.66 & 10979 & 12268.8 & 16090.7 & 20574 & 16274.2 \\
    & stdev & 265.901 & 270.652 & 382.621 & 538.007 & 823.946 & 719.359 & 763.055 \\
    \midrule
    \multirow{2}{*}{10 (402)} & means & 9791.8 & 10539.3 & 11993.9 & 13700.7 & 17737.7 & 22096.9 & 17582.3 \\
    & stdev & 893.661 & 985.953 & 1150.61 & 1206.94 & 1113.19 & 1082.1 &1240.95 \\
    \midrule 
\bottomrule
\end{tabularx}
  \caption{Pixelwise computational statistics on land cover class distribution by spectral bands of the Landsat-8 OLI/TIRS image on 2017}
  \label{tab05}
\end{table}

\begin{table}[H]
  \centering
  \begin{tabularx}{\linewidth}{p{1.3cm}p{1.0cm}p{1.0cm}p{1.0cm}p{1.0cm}p{1.0cm}p{1.0cm}p{1.0cm}p{1.0cm}}
  \toprule
  \multirow{2}{*}{Class} & \multicolumn{7}{c}{Multispectral bands of the Landsat 8 OLI/TIRS scene for \textbf{2024}} \\
    & Stat & B1 & B2 & B3 & B4 & B5 & B6 & B7 \\
    \midrule
    \multirow{2}{*}{1 (609)} & means & 7766.02 & 8030.69 & 8320.34 & 7662.62 & 7371.63 & 7425.37 & 7425.87 \\
    & stdev & 305.671 & 419.913 & 771.52 & 762.096 & 630.351 & 448.005 & 429.988 \\
    \midrule
    \multirow{2}{*}{2 (670)} & means & 7624.08 & 7809.15 & 8313.88 & 8331.51 & 14445.3 & 10705.7 & 8901.21 \\
    & stdev & 205.902 & 224.873 & 282.263 & 358.421 & 1258.89 & 839.254 & 538.858 \\
    \midrule
    \multirow{2}{*}{3 (902)} & means & 7957.84 & 8235.82 & 8901.62 & 9185.04 & 15515.2 & 12986 & 10380.8 \\
    & stdev & 246.116 & 232.907 & 266.852 & 339.025 & 1083.85 & 732.316 & 535.359 \\
    \midrule
    \multirow{2}{*}{4 (411)} & means & 8727.71 & 9082 & 9805.31 & 10220.4 & 13215.6 & 13490.9 & 12152.5 \\
    & stdev & 668.955 & 673.146 & 789.834 & 850.916 & 1438.95 & 1133.67 & 998.25 \\
    \midrule
    \multirow{2}{*}{5 (456)} & means & 8414.58 & 8770.55 & 9670.05 & 10165.1 & 17197.8 & 15038 & 11860.2 \\
    & stdev & 558.686 & 585.325 & 697.741 & 839.222 &1532.3 & 935.39 & 711.043 \\
    \midrule
    \multirow{2}{*}{6 (712)} & means & 8683.41 & 9095.66 & 10010 & 10836.8 & 15706 & 16695.9 & 13367.1 \\
    & stdev & 367.576 & 376.232 & 457.499 & 563.708 & 983.967 & 895.602 & 565.674 \\
    \midrule
    \multirow{2}{*}{7 (770)} & means & 9133.68 & 9581.99 & 10492.9 & 11575.9 & 15226.5 & 18332.2 & 15088 \\
    & stdev & 475.702 & 537.058 & 632.442 & 661.257 & 750.045 & 1015.51 & 778.522 \\
    \midrule
    \multirow{2}{*}{8 (860)} & means & 9015.93 & 9521.11 & 10624.3 & 11889.1 & 17408.1 & 19159.2 & 14827.3 \\
    & stdev & 303.977 & 345.786 & 442.905 & 672.086 & 995.622 & 796.296 & 650.634 \\
    \midrule
    \multirow{2}{*}{9 (1204)} & means & 9377.85 & 9895.28 & 10950.8 & 12359.9 & 16667.9 & 20878.6 & 16604.8 \\
    & stdev & 295.187 & 317.27 & 410.42 & 550.489 & 833.241 & 809.952 & 709.571 \\
    \midrule
    \multirow{2}{*}{10 (623)} & means & 9883.75 & 10559.2 & 11938.5 & 13760.8 & 18445.9 & 22327.5 & 17691.4 \\
    & stdev & 907.155 & 976.955 & 1174.7 & 1143.97 & 1143.34 & 1289.88 & 1481.44 \\
    \midrule 
\bottomrule
\end{tabularx}
  \caption{Pixelwise computational statistics on land cover class distribution by spectral bands of the Landsat-8 OLI/TIRS image on 2024}
  \label{tab06}
\end{table}

The overall trend in land cover type distribution by classes and the dynamics of changes is presented in Table \ref{tab07}. Here, the statistical results of the class distribution of the assigned pixels are based on the of the evaluated signature file by four years and seven multispectral bands as summarised in Table \ref{tab07} showing target land cover types for 2013--2024. Land cover class 1 (water -- Port Phillip Bay, lakes, rivers) demonstrated fluctuations with a slight increase in 2015 and 2024 which points at climate effects and higher precipitation. 

\begin{table}[H]
  \centering
  \begin{tabularx}{\textwidth}{p{0.9cm}p{0.9cm}p{0.9cm}p{0.9cm}p{0.9cm}p{0.9cm}p{0.9cm}p{0.9cm}p{0.9cm}p{0.9cm}p{0.9cm}}
  \toprule
  \multirow{2}{*}{Year} & \multicolumn{10}{c}{Classes of land cover types} \\
     & 1 & 2 &  3 & 4 & 5 & 6 & 7 & 8 & 9 & 10 \\
    \midrule
    \textbf{2013} & 1183 & 425 & 482 & 493 & 446 & 588 & 744 & 773 & 752 & 1158 \\
    \textbf{2015} & 1215 & 367 & 471 & 491 & 485 & 589 & 731 & 782 & 774 & 1100 \\
    \textbf{2017} & 1146 & 415 & 496 & 414 & 464 & 634 & 738 & 887 & 812 & 1005 \\
    \textbf{2024} & 1249 & 514 & 541 & 434 & 440 & 522 & 671 & 784 & 806 & 1256 \\
\bottomrule
\end{tabularx}
  \caption{Class distribution of the assigned pixels target land cover types for 2013--2024 during iterative classification process}
  \label{tab07}
\end{table} 

\hl{The borders of marshes and wetlands include an edge or the transitional intertidal zone between the main area of Cheetham Wetlands and the transitional zone that includes scrubs, shrubland and bushland, as well as mixed single tree cover zone. In contrast to the wetlands, the areas of farmland and cropland used for agriculture are mostly distributed at higher elevations surrounded by the wetlands and marshes in the coastal areas which are less prone to inundation due to the topographically elevated areas and drainage channels. Forest community and woody vegetation are predominantly distributed in higher elevated areas further from the beachfront area. The bushland, scrubs and shrubland class include patches of the mixed woody vegetation within wetlands and along its edges.}

The ongoing climate-environmental processes and anthropogenic activities resulted in a slight decrease of land cover classes 5 to 7 which experienced changes in occupied surface due to the cumulative effects from aridification and human pressure along the coasts of southern Australia. Land cover class 8 (built-up areas) shown the increase over the recent decade with indicates the expansion of the areas with impervious structures (e.g., such as roads, industrial construction and similar occupied lands). In contrast, the increase of pasture and grasslands (class 10) and bare areas (class 11) is notable for the recent decade which might lead to soil degradation and initiate the deforestation of the lands located in semi-arid environment, Table \ref{tab07}. 

\subsection{Interpretation of trends}
\hl{The effectiveness of different algorithms in supervised classification is based on the comparison of the information-filtering and information-extracting tools intended to anticipate  and provide relevant content regarding LULC data obtained from the processed satellite imagery. Hence, }a stable increase in land cover class 2 (bushland) is notable from 425 pixels in 2013 to 514 pixels in 2024 which signifies the decrease of forests in favour of minor plant types. The increase in land cover class 3 (farmland and cropland used for agriculture) by 2024 implies the intensification of irrigated lands used for commercial purposes. The distribution of forests (land cover class 4) demonstrated a slight decrease during the past decade from 493 pixels in 2013 to 434 pixels by 2024 which indicates a slight deforestation in the coastal areas of Victoria, Port Phillip Bay. The variations in land cover class 5 (sparse vegetation and scattered trees) indicated the initiation of the processes of restoration which follows-up the remediation of the previously abandoned lands. 

\begin{figure}[H]
  \centering
  \begin{tabular}[b]{c}
    \includegraphics[width=6.3cm]{Fig_09a.jpg} \\
    \small (a)
  \end{tabular}
  \begin{tabular}[b]{c}
    \includegraphics[width=6.3cm]{Fig_09b.jpg} \\
    \small (b) 
  \end{tabular}
    \begin{tabular}[b]{c}
    \includegraphics[width=6.3cm]{Fig_09c.jpg} \\
    \small (c)
  \end{tabular}
  \begin{tabular}[b]{c}
    \includegraphics[width=6.3cm]{Fig_09d.jpg} \\
    \small (d) 
  \end{tabular}
  \caption{Image classification on 2024 using four methods: 1) Random Forest; 2) SVM; 3) ANN-based MLPC; 4) DTClassifier. Image processing: author.}
  \label{fig09}
\end{figure}

\subsection{Comparison of approaches}

\hl{The ML algorithms that enabled to extract such information using GRASS GIS scripting methods leverage methodologies to examine satellite images based on spectral reflectance, threshold between categories for assignment of pixels, and raster cell attributes to produce maps. Here, the SVM, RF, MLPC and DTC algorithms demonstrated effectiveness in RS data processing. Using these tools, the RS data were categorised into content-based regional patterns of land cover types of southern Australia. Among these, SVM and RF scenarios are predominantly effective due to functionality of algorithms that increased accuracy results. Thus}, the SVM and RF models’ performance demonstrated similar results which is caused due to the technical parameters of these algorithms. The most distinct and accurate (over 98\% of accuracy) performance was noted by the AI algorithm of MLPC. 

The use of these methods incorporated additional information on distribution of land cover types over souther Australia, i.e., landscape dynamics, and trends over the recent decade. The overall accuracy of the SVM classifier demonstrated 97.6\%, following by 96.4\% by RF. The DTC shown lesser results yet is comparable to the tested methods with overall achieved accuracy of 96.1\%. Once the optical ML AI classifiers were determined for image processing, the time series analysis was plotted (Figure \ref{fig10}) to visualise the dynamics of the land cover types in southern Australia between 2013 to 2024 along the coastal areas of Victoria. The decrease of areas occupied by Cheetham Wetlands around Port Phillip Bay was assessed over the recent decade using RS data processing and statistical analysis.

\begin{figure}[H]
  \centering
 \begin{tabular}[b]{c}
    \includegraphics[width=6.3cm]{Fig_10a.jpg} \\
    \small (a)
  \end{tabular}
  \begin{tabular}[b]{c}
    \includegraphics[width=6.3cm]{Fig_10b.jpg} \\
    \small (b) 
  \end{tabular}
    \begin{tabular}[b]{c}
    \includegraphics[width=6.3cm]{Fig_10c.jpg} \\
    \small (c)
  \end{tabular}
  \begin{tabular}[b]{c}
    \includegraphics[width=6.3cm]{Fig_10d.jpg} \\
    \small (d) 
  \end{tabular}
  \caption{Time series of the classified images (2013-2024) processed using k-means clustering.}
  \label{fig10}
\end{figure}

%----------- subsection ----------->
\subsection{Accuracy assessment}

The accuracy assessment computed using chi-square algorithms which evaluates the probability of the pixel's assignment to the correct class is shown in Figure \ref{fig11}. Accuracy assessment was performed using the computed reject raster map which presents the threshold. The threshold \hl{in these tables }evaluates the limits according to which the pixels are assigned to the target classes correctly. It is possible using a chi square test calculated on each discriminant result at tested threshold levels of confidence. This test determines at what confidence level each pixels are classified to the correct categories of the land cover classes. Hence, the reject threshold layer contains the index to the  confidence level computed for each classified cell in the classified Landsat scene. \hl{The results of accuracy assessments evaluate the results observed image processing. }Additionally, the confusion \hl{matrices of the }separability \hl{classes} were computed to evaluate the correctness of the pixel's classification \hl{on the images }for all the years: 2013 and 2015 (Figure \ref{fig12}) and 2017 and 2024 (Figure \ref{fig13}). \hl{These figures report the class separability matrices for land cover classes which points at accuracy from this classification with the four Landsat bands. }

\hl{The Table of the Appendix A show the initial means for each band and Table in Appendix B reports the iterative cycles of image classification for each Landsat scene. The comparison of the data provide support to the overall classification procedure by GRASS GIS  in that the land cover classes are spectrally separable on the multispectral Landsat scenes and that the AI and ML algorithms of have worked quite well in these instances. The overall classification accuracy for the stratified random subset of points in 2013 was reported at 98.3\% convergency with the Kappa coefficient of 98.27\% points stable; for 2015 -- 98.3\% convergency and 98.34\% points stable; for 2017 -- convergence=98.1\% and 98.12\% points stable; finally, for 2024 we obtained 98.41\% points stable with convergence=98.4\%. Such results with the obtained Kappa coefficients above 0.80 means that the overall classification accuracy by AI /ML algorithms of GRASS GIS approach is high.}

\vspace{-10pt}
\begin{figure}[H]
  \centering
  \begin{tabular}[b]{c}
    \includegraphics[width=6.4cm]{Fig_11a.jpg} \\
    \small (a)
  \end{tabular}
  \begin{tabular}[b]{c}
    \includegraphics[width=6.4cm]{Fig_11b.jpg} \\
    \small (b) 
  \end{tabular}
    \begin{tabular}[b]{c}
    \includegraphics[width=6.4cm]{Fig_11c.jpg} \\
    \small (c)
  \end{tabular}
  \begin{tabular}[b]{c}
    \includegraphics[width=6.3cm]{Fig_11d.jpg} \\
    \small (d) 
  \end{tabular}
  \caption{Accuracy analysis of image classification (2013-2024) using chi square algorithm method.}
  \label{fig11}
\end{figure}

\begin{figure}[H]
\begin{equation}
%\setlength\arraycolsep{3pt}
\scriptsize
\begin{vmatrix}
    & 1 & 2 & 3 & 4 & 5 & 6 & 7 & 8 & 9 & 10 \\
1 & 0 &  &  &  &  &  &  &  &  &  \\
2 & 2.2 & 0 &  &  &  &  &  &  &  &  \\
3 & 3.8 & 1.2 & 0 &  &  &  &  &  &  &  \\
4 & 4.1 & 1.8 & 0.9 & 0 &  &  &  &  &  &  \\
5 & 4.1 & 1.8 & 1.0 & 0.6 & 0 &  &  &  &  &  \\                 
6 & 6.0 & 2.9 & 1.9 & 0.8 & 0.8 & 0 &  &  &  &  \\             
7 & 7.0 & 3.5 & 2.5 & 1.3 & 1.0 & 0.6 & 0 &  &  &  \\                      
8 & 7.3 & 3.8 & 2.8 & 1.6 & 1.4 & 0.9 & 0.7 & 0 &  &  \\
9 & 8.3 & 4.7 & 3.7 & 2.2 & 2.0 & 1.7 & 1.3 & 0.7 & 0 &  \\           
10 & 6.0 & 3.7 & 3.0 & 2.1 & 1.9 & 1.8 & 1.5 & 1.2 & 0.7 & 0 \\                          
\end{vmatrix}
 %
\hspace{3pt}
\begin{vmatrix}
    & 1 & 2 & 3 & 4 & 5 & 6 & 7 & 8 & 9 & 10 \\
1 & 0 &  &  &  &  &  &  &  &  &  \\
2 & 1.8 & 0 &  &  &  &  &  &  &  &  \\
3 & 3.4 & 1.1 & 0 &  &  &  &  &  &  &  \\
4 & 3.6 & 1.7 & 0.9 & 0 &  &  &  &  &  &  \\
5 & 3.9 & 2.0 & 1.1 & 0.6 & 0 &  &  &  &  &  \\  
6 & 5.3 & 2.8 & 1.9 & 0.8 & 0.9 & 0 &  &  &  &  \\   
7 & 5.8 & 3.2 & 2.4 & 1.3 & 1.2 & 0.7 & 0 &  &  &  \\ 
8 & 6.8 & 3.9 & 3.1 & 1.6 & 1.8 & 0.9 & 0.7 & 0 &  &  \\
9 & 7.7 & 4.6 & 3.9 & 2.3 & 2.4 & 1.7 & 1.2 & 0.8 & 0 &  \\
10 & 4.8 & 3.3 & 2.9 & 2.0 & 2.1 & 1.7 & 1.3 & 1.2 & 0.7 & 0 \\   
\end{vmatrix} 
\end{equation}
\caption{Class separability matrices for land cover classes: \textbf{2013} and \textbf{2015}}
\label{fig12}
\end{figure}

\begin{figure}[H]
\begin{equation}
%\setlength\arraycolsep{3pt}
\scriptsize
\begin{vmatrix}
    & 1 & 2 & 3 & 4 & 5 & 6 & 7 & 8 & 9 & 10 \\
1 & 0 &  &  &  &  &  &  &  &  &  \\
2 & 2.1 & 0 &  &  &  &  &  &  &  &  \\
3 & 3.0 & 1.4 & 0 &  &  &  &  &  &  &  \\
4 & 3.4 & 1.1 & 0.9 & 0 &  &  &  &  &  &  \\
5 & 4.3 & 2.1 & 0.7 & 1.3 & 0 &  &  &  &  &  \\    
6 & 5.2 & 2.8 & 1.0 & 2.1 & 0.8 & 0 &  &  &  &  \\  
7 & 5.7 & 3.2 & 1.4 & 2.5 & 1.0 & 0.7 & 0 &  &  &  \\ 
8 & 7.2 & 4.0 & 1.8 & 3.4 & 1.6 & 0.9 & 0.6 & 0 &  &  \\
9 & 7.6 & 4.6 & 2.3 & 4.0 & 2.2 & 1.6 & 1.0 & 0.8 & 0 &  \\
10 & 6.1 & 3.9 & 2.4 & 3.4 & 2.2 & 1.9 & 1.4 & 1.4 & 0.8 & 0 \\      
\end{vmatrix}
 %
\hspace{3pt}
\begin{vmatrix}
    & 1 & 2 & 3 & 4 & 5 & 6 & 7 & 8 & 9 & 10 \\
1 & 0 &  &  &  &  &  &  &  &  &  \\
2 & 2.0 & 0 &  &  &  &  &  &  &  &  \\
3 & 3.2 & 1.2 & 0 &  &  &  &  &  &  &  \\
4 & 2.6 & 1.6 & 1.0 & 0 &  &  &  &  &  &  \\
5 & 3.5 & 1.9 & 1.0 & 0.7 & 0 &  &  &  &  &  \\ 
6 & 4.6 & 2.9 & 1.9 & 0.9 & 0.7 & 0 &  &  &  &  \\
7 & 5.1 & 3.4 & 2.5 & 1.3 & 1.4 & 0.8 & 0 &  &  &  \\ 
8 & 5.8 & 4.0 & 3.0 & 1.6 & 1.5 & 1.0 & 0.6 & 0 &  &  \\
9 & 6.7 & 4.9 & 3.9 & 2.2 & 2.2 & 1.8 & 0.9 & 0.8 & 0 &  \\  
10 & 5.3 & 3.9 & 3.2 & 2.2 & 2.1 & 1.8 & 1.3 & 1.1 & 0.7 & 0 \\   
\end{vmatrix} 
\end{equation}
\caption{Class separability matrices for land cover classes: \textbf{2017} and \textbf{2024}}
\label{fig13}
\end{figure}

%----------- section ----------->
\section{Discussion \textcolor{blue}{and Conclusion}}

This study examined four AI applications of GRASS GIS for obtaining information on land cover changes in the coastal waters of Port Phillip Bay using data on spectral reflectance from the multi-temporal Landsat satellite data (2013-2024). The spatial extent of these changes was visualised over the study area in the surroundings of southern Victoria  State, Australia. The presented maps are applicable to identify hotspots for water eutrophication in the Port, to detect the decline of native forest ecosystems and their replacement by urban spaces and to detect landscape dynamics in a regional scale relating to environmental-climate effects and anthropogenic activities. 

The results of image classification using four different algorithms of AI and ML were compared to identify their effectiveness in GRASS GIS software \hl{and the classification algorithms of ML and AI performed well in the case of both wetland cover classes and other categories.} This approach enabled to detect the areas experiencing the most intensive land cover changes in southern Victoria, Port Phillip Bay. Pixel-level mapping enhanced the identification of regions requiring flexible methodology and investigation efficiency by focusing on specific areas such as Cheetham Wetlands which are distributed over the southern segment of the study area. \hl{A multi-disciplinary approach that integrated ML and RS data with GIS, as demonstrated in this study, ensures more insightful analysis of the interrelated processed along the coasts of southern Australia. As discussed earlier, vulnerable costal area along the Port Philipp Bay is affected both by climate-environmental factors and by anthropogenic activities. This results in highly dynamic and changing environment which was detected over the time series of the satellite images. }

\subsection{Closing remarks}
Maps play a crucial role in environmental monitoring through the visual analyses which supports environmental policy and management. Maps created using RS data within the framework of coastal environmental monitoring support decision making through the analysis of geospatial information. To support \hl{the }decision making process, the essential advantage of cartographic data visualization consists in spatio-temporal aspect of RS data and operative monitoring using time series. Using this information, we can always answer questions as "what changed in the landscapes ?" and "where do these changes take place ?" and "what is the level of changes (small to \hl{high}) ?" Hence, processing satellite images primarily \hl{supports} such decision making \hl{process }and geospatial analysis through the visualised extent of the affected areas and mapping the exact location of land cover changes. Target coastal areas affected either by the environmental-climate processes or by human activities can be highlighted to find the ways of mitigation of the environmental crisis. 

Earth observation datasets support such tasks in the environmental monitoring. This becomes possible due to the significant development of satellite missions and techniques which resulted in the increased number and variety of the available RS data. Visual forms of representation such as classified satellite images, maps, graphs and tables with obtained statistical information effectively summarise information and present it in an interpretable form which enables to highlight important zones of the endangered coastal regions on the maps.The use of such data in coastal monitoring supports high level of complexity and operability of mapping due to technical advancements in GIS and progress in programming, including the algorithms of AI and ML. Using scripting techniques of GRASS GIS enhanced through ML and AI algorithms enabled to automate and improve manual mapping, to increase the precision and quality of images and to contribute to the monitoring of coastal land use types in southern Australia. 

\subsection{Future directions}

Novel framework of coastal and marine monitoring has an multi- and interdisciplinary profile that combines diverse scientific branches: advanced RS datasets, computer graphics, data science, computing algorithms integrated with GIS methodology or visualization and spatial analysis. Such integrative approaches pave the ways to further development of marine research and open new perspectives of the environmental coastal mapping. In this paper, we demonstrated such integrated study with a case of Australian coasts and discussed practices of RS data processing in geoinformatics for environmental coastal mapping. The study was supported by an overview of modern AI and ML methods of RS data processing in modern geoinformation. Hence, the perspectives of the GIS-based integrated use of programming, ML, DL techniques and satellite image analysis were discussed and highlighted. We also provided the most essential descriptions of the novel characteristics of the AI and ML track specifically for RS data analysis. 

The challenge of future studies is to integrate and revitalise diverse sources of information for a comprehensive analysis of coastal zones. Hence, in further works that perform similar analysis in coastal areas, the integrated methods can bring research to a principally novel level. Namely, the ML-supported RS data processing by GIS can help explore coastal landscape dynamics. When applied to environmental monitoring of the coasts affected by climate hazards or erosion, such data integration promises effective tools for decision making. Hence, future studies will consider smooth integration of data and methods which is notable for the multi-disciplinary coastal monitoring supported by Earth Observation data and spatial analysis. 

To conclude, coastal monitoring is a systematic research with the goal to take decision based on the diversity of multi-format data and advanced tools of their processing. Thus, the input data can be included in various forms such as RS satellite images, cartographic data, computational tables based on ecological surveys, experimental evaluations or user opinion surveys for coastal environmental monitoring. The advantages, benefits, and usage conditions for such methods and tools for novel research directions in coastal mapping should be considered and applied in similar works. Hence, as demonstrated in this study, using modern technologies of satellite image processing, statistical GIS analysis and programming plays a crucial role in driving the environmental modelling towards the advanced level of research performed by the modern GIS. 

%----------- section ----------->
%\authorcontributions{Supervision, conceptualization, methodology, software, resources, funding acquisition, and project administration, writing---original draft preparation, methodology, software, data curation, visualization, formal analysis, validation, writing---review and editing, and investigation, P.L.}

%\funding{The publication was funded by the Editorial Office of JMSE, Multidisciplinary Digital Publishing Institute (MDPI), by providing 100\% discount for the APC of this manuscript.}

\institutionalreview{Not applicable.}

\informedconsent{Not applicable.}

\dataavailability{Not applicable.} 

\acknowledgments{The authors thank the reviewers for reading and review of this manuscript.}

\conflictsofinterest{The authors declare no conflict of interest.} 

\abbreviations{Abbreviations}{
The following abbreviations are used in this manuscript:\\

\noindent 
\begin{tabular}{@{}ll}
AI & Artificial Intelligence \\
ANN & Artificial Neural Network \\
DN & Digital Number \\
DOS & Dark Object Subtraction \\
EO & Earth Observation \\
GCP & Ground Control Points \\
GEBCO & General Bathymetric Chart of the Oceans \\
GMT & Generic Mapping Tools \\
GRASS & Geographic Resources Analysis Support System \\
GIS & Geographic Information System \\
Landsat OLI/TIRS & Landsat Operational Land Imager and Thermal Infrared Sensor \\
ML & Machine Learning \\
MLP & Multilayer Perceptron \\
\hl{MODIS} & Moderate Resolution Imaging Spectroradiometer \\
NIR & Near Infrared \\
RF & Random Forest \\
R\&D & Research and Development \\
RS & Remote Sensing \\
SPOT & Satellite Pour l'Observation de la Terre \\
SVM & Support Vector Machine \\
SWIR & Shortwave Infrared \\
USGS & United States Geological Survey \\
UTM & Universal Transverse Mercator \\
WGS84 & World Geodetic System 84 \\
WRS & Worldwide Reference System
\end{tabular}
}

%%%%%%%%%%%%%%%%%%%%%%%%%%%%%%%%%%%%%%%%%%
%% Optional
\appendixtitles{no} 
\appendixstart
\appendix

\section[\appendixname~\thesection]{Initial means for each band}\label{AppA}

\begin{table}[H]
  \centering
  \begin{tabularx}{\textwidth}{p{0.7cm}p{2.5cm}p{1.0cm}p{1.3cm}p{1.0cm}p{1.0cm}p{1.5cm}p{1.5cm}}
  \toprule
  \multirow{2}{*}{class} & \multicolumn{7}{c}{Initial mean values for multispectral bands of Landsat} \\
     & 1 - Coastal aerosol & 2 - Blue & 3 - Green & 4 - Red & 5 - NIR & 6 - SWIR-1 & 7 - SWIR-2 \\
    \midrule
    \textbf{1} & 7944.36 & 8195.91 & 8742.78 & 8861.72 & 12337.5 & 12337.4 & 10304.5 \\
    \textbf{2} & 8119.86 & 8408.03 & 9026.45 & 9267.03 & 12909.9 & 13277 & 11008.4 \\
    \textbf{3} & 8295.36 & 8620.15 & 9310.11 & 9672.33 & 13482.2 & 14216.6 & 11712.2 \\
    \textbf{4} & 8470.86 & 8832.27 & 9593.77 & 10077.6 & 14054.6 & 15156.2 & 12416.1 \\
    \textbf{5} & 8646.36 & 9044.38 & 9877.44 & 10482.9 & 14627 & 16095.8 & 13120 \\
    \textbf{6} & 8821.86 & 9256.5 & 10161.1 & 10888.3 & 15199.4 &17035.4 & 13823.8  \\
    \textbf{7} & 8997.36 & 9468.62 & 10444.8 & 11293.6 & 15771.7 & 17975 & 14527.7 \\
    \textbf{8} & 9172.86 & 9680.74 & 10728.4 & 11698.9 & 16344.1 & 18914.6 & 15231.5 \\
    \textbf{9} & 9348.36 & 9892.86 & 11012.1 & 12104.2 & 16916.5 & 19854.2 & 15935.4 \\
    \textbf{10} & 9523.86 & 10105 & 11295.8 & 12509.5 & 17488.8 & 20793.8 & 16639.3 \\
\bottomrule
\end{tabularx}
  \caption{Initial mean values for pixels in each band assigned to target land cover types for 2013 before the iterative classification process}
  \label{tabA1} 
\end{table} 

\begin{table}[H]
  \centering
  \begin{tabularx}{\textwidth}{p{0.7cm}p{2.5cm}p{1.0cm}p{1.3cm}p{1.0cm}p{1.0cm}p{1.5cm}p{1.5cm}}
  \toprule
  \multirow{2}{*}{class} & \multicolumn{7}{c}{Initial mean values for multispectral bands of Landsat} \\
     & 1 - Coastal aerosol & 2 - Blue & 3 - Green & 4 - Red & 5 - NIR & 6 - SWIR-1 & 7 - SWIR-2 \\
    \midrule
    \textbf{1} & 7964.53 & 8167.16 & 8635.84 & 8723.91 & 12003.5 & 11922.3 & 10077.8 \\
    \textbf{2} & 8143.64 & 8371.91 & 8901.87 & 9095.84 & 12597.6 & 12891.5 & 10824.4  \\
    \textbf{3} & 8322.74 & 8576.66 & 9167.9 & 9467.76 & 13191.7 & 13860.6 & 11571 \\
    \textbf{4} & 8501.84 & 8781.41 & 9433.92 & 9839.69 & 13785.8 & 14829.8 & 12317.6 \\
    \textbf{5} & 8680.94 & 8986.16 & 9699.95 & 10211.6 & 14379.8 & 15798.9 & 13064.2 \\
    \textbf{6} & 8860.04 & 9190.91 & 9965.98 & 10583.5 & 14973.9 & 16768.1 & 13810.8 \\
    \textbf{7} & 9039.15 & 9395.66 & 10232 & 10955.5 & 15568 & 17737.2 & 14557.5 \\
    \textbf{8} & 9218.25 & 9600.41 & 10498 & 11327.4 & 16162.1 & 18706.4 & 15304.1 \\
    \textbf{9} & 9397.35 & 9805.16 & 10764.1 & 11699.3 & 16756.2 & 19675.5 & 16050.7 \\
    \textbf{10} & 9576.45 & 10009.9 & 11030.1 & 12071.2 &17350.2 & 20644.7 & 16797.3 \\
\bottomrule
\end{tabularx}
  \caption{Initial mean values for pixels in each band assigned to target land cover types for 2015 before the iterative classification process}
  \label{tabA2} 
\end{table} 

\begin{table}[H]
  \centering
  \begin{tabularx}{\textwidth}{p{0.7cm}p{2.5cm}p{1.0cm}p{1.3cm}p{1.0cm}p{1.0cm}p{1.5cm}p{1.5cm}}
  \toprule
  \multirow{2}{*}{class} & \multicolumn{7}{c}{Initial mean values for multispectral bands of Landsat} \\
     & 1 - Coastal aerosol & 2 - Blue & 3 - Green & 4 - Red & 5 - NIR & 6 - SWIR-1 & 7 - SWIR-2 \\
    \midrule
        \textbf{1} & 7956.33 & 8207.65 & 8763.65 & 8860.93 & 12306 & 11946.7 & 10001.2 \\
    \textbf{2} & 8117.93 & 8401.46 & 9023.47 & 9249.65 & 12895.4 & 12888 & 10677.9 \\
    \textbf{3} & 8279.53 & 8595.27 & 9283.28 & 9638.37 & 13484.7 & 13829.4 & 11354.6 \\
    \textbf{4} & 8441.13 & 8789.08 & 9543.09 & 10027.1 & 14074.1 & 14770.7 & 12031.3 \\
    \textbf{5} & 8602.73 & 8982.89 & 9802.91 & 10415.8 & 14663.4 & 15712 & 12708 \\
    \textbf{6} & 8764.32 & 9176.7 & 10062.7 & 10804.5 & 15252.8 & 16653.4 & 13384.7 \\
    \textbf{7} & 8925.92 & 9370.51 & 10322.5 & 11193.3 & 15842.1 & 17594.7 & 14061.4 \\
    \textbf{8} & 9087.52 & 9564.32 & 10582.3 & 11582 & 16431.5 & 18536.1 & 14738.1 \\
    \textbf{9} & 9249.12 & 9758.14 & 10842.2 & 11970.7 & 17020.8 & 19477.4 & 15414.8 \\
    \textbf{10} & 9410.72 & 9951.95 & 11102 & 12359.4 & 17610.2 & 20418.7 & 16091.6 \\
\bottomrule
\end{tabularx}
  \caption{Initial mean values for pixels in each band assigned to target land cover types for 2017 before the iterative classification process}
  \label{tabA3} 
\end{table} 

\begin{table}[H]
  \centering
  \begin{tabularx}{\textwidth}{p{0.7cm}p{2.5cm}p{1.0cm}p{1.3cm}p{1.0cm}p{1.0cm}p{1.5cm}p{1.5cm}}
  \toprule
  \multirow{2}{*}{class} & \multicolumn{7}{c}{Initial mean values for multispectral bands of Landsat} \\
     & 1 - Coastal aerosol & 2 - Blue & 3 - Green & 4 - Red & 5 - NIR & 6 - SWIR-1 & 7 - SWIR-2 \\
    \midrule
        \textbf{1} & 7883.17 & 8161.79 & 8731.22 & 8868.39 & 12420.6 & 11749.3 & 9919.64 \\
    \textbf{2} & 8066.83 & 8374.7 & 9010.16 & 9289.44 & 13074.4 & 12751.4 & 10645.2 \\
    \textbf{3} & 8250.49 & 8587.61 & 9289.1 & 9710.5 & 13728.3 & 13753.5 & 11370.7 \\
    \textbf{4} & 8434.15 & 8800.53 & 9568.04 & 10131.5 & 14382.2 & 14755.5 & 12096.3 \\
    \textbf{5} & 8617.81 & 9013.44 & 9846.97 & 10552.6 & 15036 & 15757.6 & 12821.8 \\
    \textbf{6} & 8801.47 & 9226.35 & 10125.9 & 10973.7 & 15689.9 & 16759.6 & 13547.4 \\
    \textbf{7} & 8985.13 & 9439.26 & 10404.9 & 11394.7 & 16343.7 & 17761.7 & 14272.9 \\
    \textbf{8} & 9168.79 & 9652.17 & 10683.8 & 11815.8 & 16997.6 & 18763.8 & 14998.4 \\
    \textbf{9} & 9352.45 & 9865.08 & 10962.7 & 12236.8 & 17651.4 & 19765.8 & 15724 \\
    \textbf{10} & 9536.11 & 10078 & 11241.7 & 12657.9 & 18305.3 & 20767.9 & 16449.5 \\
\bottomrule
\end{tabularx}
  \caption{Initial mean values for pixels in each band assigned to target land cover types for 2024 before the iterative classification process}
  \label{tabA4} 
\end{table} 

\section[\appendixname~\thesection]{Iterative cycles of image classification}\label{AppB}

\begin{table}[H]
  \centering
  \begin{tabularx}{\textwidth}{p{1.6cm}p{1.8cm}p{0.6cm}p{0.6cm}p{0.6cm}p{0.6cm}p{0.6cm}p{0.6cm}p{0.6cm}p{0.6cm}p{0.6cm}p{0.6cm}}
  \toprule
  \multirow{2}{*}{2013} & \multicolumn{10}{c}{Pixels' distribution by categories of land cover classes} \\
     & points stable  & 1 & 2 &  3 & 4 & 5 & 6 & 7 & 8 & 9 & 10 \\
    \midrule
    Iteration 1 & 79.91\% & 747 & 855 & 508 & 460 & 478 & 595 & 688 & 802 & 1014 & 897 \\
    Iteration 2 & 88.69\% & 558 & 897 & 671 & 456 & 477& 594 &683 & 861& 1075 &772 \\
    Iteration 3 & 90.81\% & 513 & 801 & 806 & 491 & 449 & 603 & 693 & 893 & 1108 & 687 \\
    Iteration 4 & 93.02\% & 510  & 733 & 867 & 506 & 448 & 618 & 701 & 910 & 1129 & 622 \\
    Iteration 5 & 95.30\% & 508 & 690 & 912 & 515 & 417 & 641  & 727& 910 &1145 &579\\
    Iteration 6 & 96.39\% & 507 & 671 & 924 &531 &398 & 668 & 698 & 948 &1149  &550 \\
    Iteration 7 & 97.30\% &  507 & 655 & 924 & 556 & 387& 701& 658 & 966 &1163 & 527\\
    Iteration 8 & 97.90\% &  507 & 646 & 905 &585 &381 & 734 & 613 &987 & 1175 &511 \\
    Iteration 9 & 98.27\% & 507 & 638  & 887 &617 & 378 &748 & 584 &1004 &1175  &506\\
\bottomrule
\end{tabularx}
  \caption{Iterative process of pixels' assignment to target classes during image classification: 2013}
  \label{tabB1}
\end{table}

\begin{table}[H]
  \centering
  \begin{tabularx}{\textwidth}{p{1.6cm}p{1.8cm}p{0.6cm}p{0.6cm}p{0.6cm}p{0.6cm}p{0.6cm}p{0.6cm}p{0.6cm}p{0.6cm}p{0.6cm}p{0.6cm}}
  \toprule
  \multirow{2}{*}{2015} & \multicolumn{10}{c}{Pixels' distribution by categories of land cover classes} \\
      & points stable & 1 & 2 &  3 & 4 & 5 & 6 & 7 & 8 & 9 & 10 \\
    \midrule
    Iteration 1 & 77.04\% & 766  & 831  &451& 500  &490 &594  &722 &763 &1028 & 860 \\
    Iteration 2 & 88.08\% & 627 & 815 &610 & 509 &470& 609 &708 &848 &1094&715 \\
    Iteration 3 & 89.31\% & 584 & 751& 685 & 567 & 474& 591 &711 &912&1122 &608 \\
    Iteration 4 & 91.53\% & 575  & 693&760&585 & 440& 629 & 717 & 939& 1144 &523 \\
    Iteration 5 & 93.52\% & 573 & 663 & 791 & 612  &396 & 667 & 707 & 979&1169 &448 \\
    Iteration 6 & 96.20\% & 573 & 646 & 787 & 636 &383 &719 &664 &994 & 1203 & 400 \\
    Iteration 7 & 97.36\% & 571 & 636 &776 & 666 &385& 750 & 633  &1006 &1223 &359 \\
    Iteration 8 & 97.79\% & 569 & 622 &779 & 688 &386 &776 &603&1021&1236&325\\
    Iteration 9 & 98.34\% & 568& 615 &781&703  &391&796&582& 1026&1245&298 \\
\bottomrule
\end{tabularx}
  \caption{Iterative process of pixels' assignment to target classes during image classification: 2015}
  \label{tabB2}
\end{table}

\begin{table}[H]
  \centering
  \begin{tabularx}{\textwidth}{p{1.6cm}p{1.8cm}p{0.7cm}p{0.6cm}p{0.6cm}p{0.6cm}p{0.6cm}p{0.6cm}p{0.6cm}p{0.6cm}p{0.6cm}p{0.6cm}}
  \toprule
  \multirow{2}{*}{2017} & \multicolumn{10}{c}{Pixels' distribution by categories of land cover classes} \\
    & points stable & 1 & 2 &  3 & 4 & 5 & 6 & 7 & 8 & 9 & 10 \\
    \midrule
    Iteration 1 & 74.38\% & 695 & 882& 485 &429 &448& 618 & 766 & 849& 1079& 760 \\
    Iteration 2 & 91.54\% & 602 & 844 & 621 & 429 & 456 & 599 & 775 & 930 & 1100 &  655 \\
    Iteration 3 & 91.93\% & 585 & 772 & 687 & 445 &479 & 598 & 780 & 978 & 1094 & 593 \\
    Iteration 4 & 93.21\% & 583 & 737 & 704 & 447 & 500 & 588 & 809 & 1006 & 1093 &  544\\
    Iteration 5 & 95.31\% & 583 & 723 & 707 & 435 & 528 & 588 & 818 & 1049 & 1071 & 509 \\
    Iteration 6 & 96.38\% & 583 &724 & 699 & 429 & 534 & 601 & 813 &1086 &1061 & 481 \\
    Iteration 7 & 96.41\% & 583 & 722 & 688 & 433 & 531 & 622 & 806 & 1106 & 1066 & 454 \\
    Iteration 8 & 96.36\% & 583 &722 & 669 & 446 & 533 & 656& 785 & 1105 & 1083 &429 \\
    Iteration 9 & 96.65\% & 583 & 723 &643 & 468 & 531 & 691 & 771 &1089 & 1094 & 418 \\
    Iteration 10 & 96.82\% & 582 & 726 & 621 &486 &538 &723 &739 &1074  & 1113 & 409 \\
    Iteration 11 & 97.10\% & 582 & 724 & 586 &524  & 537 &757 & 708 &1060 &1129 &404 \\
     Iteration 12 & 97.39\% & 582 & 704 & 552 & 582 & 536 & 786 &682 &1047 &1139 &401 \\
     Iteration 13 & 97.62\% &  581 &683  &523 &633 &536 &816 & 652 & 1044 & 1142 &401 \\
     Iteration 14 & 98.12\% & 581 & 673 & 494 &674 &534 & 836 & 629 &1048 & 1140 &402 \\
\bottomrule
\end{tabularx}
  \caption{Iterative process of pixels' assignment to target classes during image classification: 2017}
  \label{tabB3}
\end{table}

\begin{table}[H]
  \centering
  \begin{tabularx}{\textwidth}{p{1.4cm}p{1.8cm}p{0.7cm}p{0.6cm}p{0.6cm}p{0.6cm}p{0.6cm}p{0.6cm}p{0.6cm}p{0.6cm}p{0.6cm}p{0.6cm}}
  \toprule
  \multirow{2}{*}{2024} & \multicolumn{10}{c}{Pixels' distribution by categories of land cover classes} \\
    & points stable  & 1 & 2 &  3 & 4 & 5 & 6 & 7 & 8 & 9 & 10 \\
    \midrule
    Iteration 1 & 78.77\% & 763 &1016&519 &447 & 436 & 517 & 673 &772 & 1086 & 988 \\
    Iteration 2 & 89.39\% & 637 & 969 & 688 & 458 &430 &523 & 665 &845 &1148 & 854 \\
    Iteration 3 & 90.56\% & 616 & 845 & 794 & 491 &  458 & 493 & 685  & 913 & 1160 & 762 \\
    Iteration 4 & 93.17\% & 612 & 772 & 854 &465 &477 &523 & 686 &967 & 1160 &701 \\
    Iteration 5 & 95.08\% & 611 &  734 & 885 &  436 &  477 & 561 & 688 &  1000& 1161 &  664 \\
    Iteration 6 & 96.16\% & 609 &706 & 897 &429 &469 &598 &702 & 989 &1173 &645 \\
    Iteration 7 & 97.06\% & 609 & 685 & 907 &421& 460 & 646 &714 & 952 &1193 &630 \\
    Iteration 8 & 97.73\% & 609  &675  &907 &415 & 452  &687 &739  & 910  &1197  & 626 \\
    Iteration 9 & 98.41\% & 609& 670 & 902&411 &456 &712 & 770 & 860 &1204 &623 \\
\bottomrule
\end{tabularx}
  \caption{Iterative process of pixels' assignment to target classes during image classification: 2024}
  \label{tabB4}
\end{table}

%%%%%%%%%%%%%%%%%%%%%%%%%%%%%%%%%%%%%%%%%%
\begin{adjustwidth}{-\extralength}{0cm}
%\printendnotes[custom] % Un-comment to print a list of endnotes

\reftitle{References}
\nocite{*}

%=====================================
% References, variant A: external bibliography
%=====================================
\bibliography{Australia.bib}

%%%%%%%%%%%%%%%%%%%%%%%%%%%%%%%%%%%%%%%%%%
\end{adjustwidth}
\end{document}
